\section{Diskussion}
\label{sec:Diskussion}

\subsection{Drehschieberpumpe}
Die ermittelten Saugvermögen bei der Evakuierungsmethode für die Drehschieberpumpe
der Firma ILMVAC Typ 300883/AKD16 beliefen sich auf 
\begin{align*}
    S_\text{eff1}^\text{eva} =& \qty{1.36+-0.14}{\liter\per\second}, \\
    S_\text{eff2}^\text{eva} =& \qty{0.48+-0.05}{\liter\per\second}, \\
    S_\text{eff3}^\text{eva} =& \qty{0.19+-0.2}{\liter\per\second}.
\end{align*}
Dabei ist $S_\text{eff1}^\text{eva}$ das Saugvermögen im Druckbereich von 0 bis 
$\qty{150}{\milli\bar}$, $S_\text{eff2}^\text{eva}$ in einem Bereich von 150 bis 
$\qty{290}{\milli\bar}$ und $S_\text{eff3}^\text{eva}$ im Bereich von 290 bis 
knapp $\qty{290}{\milli\bar}$. Laut Hersteller liegt der Wert des Saugvermögens 
in einem Bereich von $\qty{4.6}{m^3/h}$ bis $\qty{5.5}{m^3/h}$ bzw. 
$\qty{1.28}{\liter\per\second}$ bis $\qty{1.53}{\liter\per\second}$. Der Wert für 
$S_\text{eff1}^\text{eva}$ liegt exakt in diesem Intervall und stimmt innerhalb
der Messunsicherheit von $±\qty{0.5}{m^3/h}$ mit der Herstellerangabe überein.
Dies ist plausibel, da im hohen Druckbereich (Grobvakuum, etwa oberhalb von
$\qty{1}{mbar}$) die laminare Strömung dominiert und die Pumpe nahe ihrem
Optimalwert arbeitet. Rückströmungseffekte und Desorption sind bei diesen Drücken
noch vernachlässigbar.
% \vspace{0.5em}
\\
Die deutlich niedrigeren Werte $S_\text{eff2}^\text{eva}$ und 
$S_\text{eff3}^\text{eva}$ sind keine Messfehler, sondern eine direkte Folge der 
Druckabhängigkeit des Saugvermögens. Mit sinkendem Druck treten einige Effekte
zunehmend in den Vordergrund. Dazu gehört der Einfluss der Molekularströmung: 
die mittlere freie Weglänge wird größer als die Abmessungen der Pumpe, was zur 
Folge hat, dass der gerichtete Gastransport ineffizienter wird. Darüber hinaus 
kommt es zu Rückströmungseffekten, wie bereits in \autoref{sec:dsp} angesprochen.
Durch Spalte zwischen Rotor und Gehäuse oder Ventilen diffundieren Gasmoleküle 
zurück in den Rezipienten. Da die Drehschieberpumpe mit Öl als Schmiermittel 
arbeitet ist es auch nicht ausgeschlossen, dass Teilchen von diesem in den 
Rezipienten übergehen. Dazu kommen auch Wasserteilchen, welche sich an den 
Wänden abgesetzt haben, was auf Van-der-Waals-Kräfte oder Wasserstoffbrücken 
an der Metalloberfläche zurückzuführen ist. Folglich ist die sinkende Steigung, 
wie in \autoref{fig:10} zu sehen, genau das erwartete Verhalten. Die Herstellerangaben
beziehen sich auf den optimalen Arbeitsbereich bei Drücken oberhalb von $\qty{1}{mbar}$.
\vspace{0.5em}
\\
Zusätzlich bestätigen die Leckratenmessungen diesen Befund. Gemessen wurden:
\begin{align*}
    S_{\text{eff}, 0.2\,\text{mbar}}^{\text{leck}}  &= \qty{0.48(15)}{\liter\per\second} \\
    S_{\text{eff}, 6.3\,\text{mbar}}^{\text{leck}}  &= \qty{1.59(50)}{\liter\per\second} \\
    S_{\text{eff}, 56.4\,\text{mbar}}^{\text{leck}} &= \qty{1.24(39)}{\liter\per\second} \\
    S_{\text{eff}, 86\,\text{mbar}}^{\text{leck}}   &= \qty{1.53(48)}{\liter\per\second}
\end{align*}
Bei $\qty{0.2}{mbar}$ ergibt sich $S_\text{eff}^\text{leck} = \qty{0.48(15)}{\liter\per\second}$, was in
Übereinstimmung mit $S_\text{eff2}^\text{eva}$ steht, während bei etwa 
$\qty{56.4}{mbar}$ ein Vermögen von $\qty{1.24+-0.39}{m^3/h}$ gemessen werden, 
kompatibel mit $S_\text{eff1}^\text{eva}$. Die beiden unabhängigen Methoden 
liefern also konsistente Ergebnisse.
\noindent Insgesamt sind die ermittelten Werte physikalisch legitim und zeigen,
dass das Saugvermögen einer Drehschieberpumpe kein individueller Materialparameter
ist, sondern stark vom Druckbereich abhängt.

\subsection{Turbomolekularpumpe}
Für die Turbomolekularpumpe SST81 der Firma ILMVAC (betrieben bei $\qty{1350}{\hertz}$)
ergab sich das effektive Saugvermögen bei der Evakuierungsmessung als
\begin{equation*}
    S_{\text{eff}, 1}^\text{Eva} = \qty{15.4(16)}{\liter\per\second}
\end{equation*}
in dem Bereich der ersten 22 Sekunden und 
\begin{equation*}
    S_{\text{eff}, 2}^\text{Eva} = \qty{0.41(5)}{\liter\per\second}.
\end{equation*}
für den späteren Verlauf.
Laut Hersteller sollte ein Vermögen von $\qty{77}{\liter\per\second}$ erreicht 
werden können, solange Stickstoff mit der Pumpe gefördert werden soll. In diesem 
Experiment wird mit Luft gearbeitet, diese besteht nur etwa zu 80 \% aus Stickstoff, 
ein exakter Vergleich ist also ohne Weiteres nicht möglich. In Anbetracht dessen 
liegt die Abweichung bei mindestens 80 \%. Diese signifikante Abweichung resultiert 
aus einigen Faktoren, wie etwa dem Aspekt, dass das effektive Saugvermögen bestimmt 
wurde. Das bedeutet, dass die Berechnung keine leitwertbehafteten Verbindugnen 
wie Rohre und Ventile berücksichtigt. Da Messungen im Bereich der Größenordnung 
von $\qty{e-5}{\milli\bar}$ liegen, herrscht molekulare Strömung. Hier ist der 
Einfluss des Leitwerts besonders stark, da das Nennsaugvermögen der Turbopumpe
sehr hoch ist, während der Leitwert typischer Rohrleitungen vergleichsweise
gering ausfällt. Weiter könnte die Rohrvergengung zwischen Turbopumpe und dem 
Rezipienten gemessene Abweichungen erklären. Diese Verengung hätte mit einem 
verringerten Leitwert den Effekt, im molekularen Strömungsbereich das effektive 
Saugvermögen drastisch zu reduzieren (unabhängig vom Nennsaugvermögen der Pumpe).
Dadurch erklärt sich ein Großteil der gemessenen Abweichung, ohne dass ein Defekt 
der Pumpe oder ein Messfehler angenommen werden muss (vergleiche \autoref{fig:7}).
\vspace{0.5em}
\\
Etwas bessere Ergebnisse liefern die Messungen mittels Leckratenmethode, welche bei 
einem Druck von $p_g = \qty{56.96e-5}{\milli\bar}$ fast ein Saugvermögen 
von mehr als $\qty{40}{\liter\per\second}$ ergaben. Insgesamt liegt die Abweichung 
noch bei knapp 50.9 bis 73.9\%. Ein Großteil der Abweichungen ist auf die
Messgenauigkeit der verwendeten PKR-360 Sensoren zurückzuführen, die vom Hersteller 
im Bereich unterhalb von $\qty{100}{\milli\bar}$ mit $30\,\%$ vom Messwert
angegeben ist.
\noindent Zusammenfassend sind die ermittelten Werte physikalisch gerechtfertig,
wenn die realen experimentellen Randbedingungen beachtet werden. Die Abweichung
vom Herstellerwert ist kein Widerspruch zur zugrundeliegenden Physik, sondern
eine direkte Konsequenz des realen Aufbaus, der nicht die idealen 
Vergleichsbedingungen eines Normtests erfüllt.

\subsection{Engstelle und Leitwert}
Der Einbau einer Engstelle in Form eines Testrohrs führt zu einer drastischen 
Reduktion des effektiven Saugvermögens von knapp $\qty{35.1}{\liter\per\second}$
auf $\qty{2,68}{\liter\per\second}$. Die gemessene Druckdifferenz zwischen den
beiden PKR-360 Sensoren bestätigt diesen Effekt: Vor der Engstelle liegt ein
niedrigerer Druck an als dahinter. Dies zeigt, dass bereits ein kurzes Rohr mit
kleinem Querschnitt einen dominanten Strömungswiderstand darstellt. Die Pumpe
kann ihr Nennsaugvermögen daher nur dann entfalten, wenn die Verbindungsleitungen
ausreichend dimensioniert sind. Eine demenstprechende Abweichung beider Werte
ist physikalisch plausibel und unterstreicht die Relevanz des Leitwerts.

\subsection{Störeinflüsse durch Konvektion und Kondensation}
\subsubsection{Konvektion}
Ein systematischer Fehler, welcher bei der Interpretation der Fehler berücksichtigt 
werden muss, liegt bei der Konvektion. Bei der Leckratenmessung wird die
Druckanstiegsrate $\frac{\increment p}{\increment t}$ gemessen und in ein 
effektives Saugvermögen umgerechnet. Das Pirani-Messprinzip, welches sowohl im 
im TPG-202 als auch im PKR-360 verwendet wird ist jedoch konvektionsempfindlich
\cite{glossar}. Temperaturunterschiede im Rezipienten führen zu Gasströmungen,
die zusätzliche Wärme vom Messdraht abführen und somit einen höheren Druck
anzeigen, als tatsächlich vorhanden ist.

\subsubsection{Kondensation als thermodynamische Pumpe im Rezipienten}
Während der Evakuierung kann es dazu kommen, dass Wasserdampf an kalten Stellen
des Rezipienten kondensiert. Dieser Phasenübergang reduziert die Teilchenzahl 
in der Gasphase, ohne dass die Pumpe tatsächlich aktiv ist. Es wirkt wie eine
zusätzliche, nicht kontrollierbare Pumpe. Im vorliegenden Experiment ist dieser 
Effekt jedoch vernachlässigbar klein, da die Messzeiten kurz und die 
Desorptionsraten gering sind, was sich darin äußert, dass bei der Evakuierungskurve 
das Saugvermögen wesentlich kleiner ist, als durch die Literatur vorgegeben. Ein 
Einfluss dieses Effekts würde sich durch eine Überschätzung des Saugvermögens 
bemerkbar machen, nicht aber durch einen zu geringen Wert.

\subsection*{Hinweis zum Einsatz von KI-Tools}
Teile der Diskussion wurden unter Verwendung des KI-Modells DeepSeek
erstellt und anschließend überprüft und angepasst.